%! Introducing Statistics: What Can We Learn from Data? — Unit 1, Lesson 1.0 — Lesson Plan
% EFFL plan (experience first, formalize later). Teacher-facing — the spoiler rule does NOT
% apply here: use the formal AP vocabulary freely. Every box is `skillbox` (breakable) with a
% \boxguard ahead of it; this course has no `fixedskillbox` and no tier boxes.
\documentclass[10pt]{article}
\usepackage{apstats-article}
\usepackage{apstats-boxes}

\newcommand{\UnitNumberName}{Unit 1: Exploring One-Variable Data \quad}
\newcommand{\LessonNumberName}{Lesson 1.0: Introducing Statistics: What Can We Learn from Data?}

\begin{document}
\begin{center}
    {\Large\bfseries \CourseName \\
    {\small\bfseries \UnitNumberName \LessonNumberName}}
\end{center}

\begin{tcolorbox}[colback=sky, colframe=navy, boxrule=0.9pt, arc=2mm,
    left=2mm, right=2mm, top=1.5mm, bottom=1.5mm]
    \textbf{Primary Objective:} Students will explain that a number carries meaning only when
    placed in context, identify the individuals and variables in a data set, and distinguish a
    statistical question --- one that anticipates variability and is answered by collecting
    data --- from a question with a single definite answer.
    \\[2pt]
    \textbf{CED alignment:} Topic 1.1 \quad$\cdot$\quad Big Idea: VAR (Variation \& Distribution)
    \quad$\cdot$\quad Skill 1.A \quad$\cdot$\quad LO VAR-1.A \quad$\cdot$\quad EK VAR-1.A.1
    \\[2pt]
    \textbf{Lesson model:} \emph{Experience first, formalize later.} Students work the activity
    in small groups using only prior knowledge; the teacher circulates with questions, cues, and
    prompts (not answers); the debrief attaches the formal vocabulary to what the groups already
    found; the application and practice apply it to new contexts. No separate direct-instruction
    block, and \textbf{no vocabulary before the debrief.}
\end{tcolorbox}

\renewcommand{\arraystretch}{1.4}
\boxguard
\begin{skillbox}[Learning Targets \& Key Understandings]{goldbox}
\begin{tabularx}{\linewidth}{@{}X|X@{}}
\textbf{Learning targets (student language)} & \textbf{Key understandings (the ``why'')} \\[2pt]
\hline
\vspace{-6pt}
\begin{itemize}[leftmargin=1.1em, nosep, after=\vspace{2pt}]
  \item I can explain why a bare number tells me nothing, and name the context it needs.
  \item I can identify the individuals and the variables in a data set, and tell a variable
        apart from one of its values.
  \item I can decide whether a question is statistical, and justify it by pointing to what varies.
  \item I can write a statistical question of my own about a given data set.
\end{itemize}
&
\vspace{-6pt}
\begin{itemize}[leftmargin=1.1em, nosep, after=\vspace{2pt}]
  \item Numbers convey meaningful information only when placed in context (EK VAR-1.A.1).
  \item Variability is the reason statistics exists: if every individual gave the same value,
        one measurement would answer everything.
  \item \textbf{Target misconception:} that any question with a numerical answer, or any
        question needing the whole table, is statistical. \emph{``How many dogs are on the
        board?''} needs every row and returns a number --- and is \textbf{not} statistical,
        because its answer does not differ from one dog to the next. The discriminator is
        \emph{variation across individuals}, not arithmetic.
\end{itemize}
\end{tabularx}
\end{skillbox}

\boxguard
\begin{skillbox}[{Vocabulary, Concepts, \& Theorems --- attached in the debrief, not before}]{greenbox}
\begin{minipage}[t]{0.48\linewidth}
    \begin{itemize}
        \item \textbf{Statistics} --- the science of collecting, organizing, analyzing, and
              drawing conclusions from data
        \item \textbf{Data} --- recorded values, together with the context that gives them meaning
        \item \textbf{Context} --- who/what was measured, when, where, and in what units
              (EK VAR-1.A.1)
        \item \textbf{Individual} --- the person, animal, or object described by one row
    \end{itemize}
\end{minipage}
\hfill
\begin{minipage}[t]{0.48\linewidth}
    \begin{itemize}
        \item \textbf{Variable} --- a characteristic recorded for each individual; a column
              heading, never one of the entries beneath it
        \item \textbf{Value} --- one entry under a variable
        \item \textbf{Variability} --- values differ across individuals; why one measurement
              is not enough
        \item \textbf{Statistical question} --- anticipates variability and is answered by
              collecting data
    \end{itemize}
\end{minipage}
\end{skillbox}

% A tabularx never splits, so guard generously ahead of it.
\boxguard[30]
\begin{skillbox}[Lesson at a Glance --- \MeetingLength]{sky}
\renewcommand{\arraystretch}{1.25}
\begin{tabularx}{\linewidth}{@{}l c X X@{}}
\textbf{Phase} & \textbf{Min} & \textbf{Students} & \textbf{Teacher} \\
\hline
Warm-Up & 5 & Seven bare numbers; three questions, alone then with a neighbor & Circulate; debrief item~3 aloud and leave it on the board all period \\
Experience \& Formalize: Activity & 20 & \emph{The Adoption Board}, scenarios 1--2, groups of 3--4 & Circulate; questions, cues, prompts --- \emph{not answers}; choose group work to display \\
Debrief: Formalize & 12 & Share findings; fill QuickNotes as each name lands & Name each idea on the groups' own work, in a second color; fill QuickNotes together \\
Application & 6 & \emph{The Gym Sign-Up}, worked together & Lead it; students hold the pen, you hold the questions \\
Check Your Understanding & 8 & 5 exercises, pairs $\to$ solo (\emph{not scored}) & Circulate; spot-check item~5, do not collect for a grade \\
Close \& Assign & 4 & Record the homework; hear the preview & Sort item-5 responses to decide how Lesson~1.1 opens \\
\end{tabularx}
\end{skillbox}

\boxguard[20]
\begin{skillbox}[Warm-Up --- Activate Prior Knowledge]{sky}
\begin{minipage}[t]{0.48\linewidth}
    \textbf{The three items and what each seeds}
    \begin{itemize}
        \item \textbf{1.} Seven numbers, no context: \emph{``What do these tell you?''} The
              intended answer is \emph{``I can't tell.''} \textbf{Seeds:} context / EK VAR-1.A.1.
        \item \textbf{2.} ``Three things you'd need to know.'' \textbf{Seeds:} the W questions ---
              their own list becomes the QuickNotes bullet.
        \item \textbf{3.} ``How tall is Maria?'' vs.\ ``How tall are the students in this class?''
              \textbf{Seeds:} the statistical-question test.
    \end{itemize}
\end{minipage}
\hfill
\begin{minipage}[t]{0.48\linewidth}
    \textbf{Running it}
    \begin{itemize}
        \item \textbf{Resist rescuing them on item~1.} The discomfort is the lesson. Three
              minutes, then ask only \emph{``what would you need to know?''}
        \item \textbf{Name nothing.} Do not say \emph{context}, \emph{variable}, or
              \emph{statistical question} in the warm-up debrief --- those arrive after the
              activity. Write their own phrasings on the board and leave them there.
        \item Debrief item~3 aloud; it reappears as the crux in scenario~2.
        \item \emph{This is the course's first lesson --- there is no prior unit to spiral.}
    \end{itemize}
\end{minipage}
\end{skillbox}

\boxguard
\begin{skillbox}[Experience \& Formalize --- The Activity: \emph{The Adoption Board}]{sky}
\textbf{Launch (2 min).} Project the board. \emph{``This is one shelter's adoption board. Visitors
read it, and so does the director --- but not for the same things. Work the two problems with
your group using only what you already know. For every answer, be ready to show me
\textbf{where on the board} you see it.''} Groups of 3--4. Do not preview any vocabulary.

\begin{multicols}{2}
\textbf{What students do.} \emph{Scenario 1} builds the whole toolkit on one table: what a row
stands for, counting rows and columns, heading-vs-entry (1b), the look-it-up / read-the-column
contrast (1c), and why a column of identical entries would be worthless (1d --- variability,
unnamed).

\emph{Scenario 2} is the contrast case. Item 2a sorts four director's questions into ``one spot''
or ``many rows.'' \textbf{2b is the crux:} two questions both need many rows, but only one has an
answer that changes from dog to dog. Items 2c--2d turn the flyer claim into the
summary-does-not-predict-an-individual idea.

\columnbreak
\textbf{What the teacher does.} Circulate. Listen before speaking. When a group stalls, give a
\emph{question, cue, or prompt}, not an answer:
\begin{itemize}[leftmargin=1.1em, nosep]
  \item \textbf{1b:} ``Point at where you see the word `Breed.' Now point at `Beagle.' Are those
        the same kind of thing?''
  \item \textbf{1d:} ``Cover the Days column with your hand. What did you just lose?''
  \item \textbf{2b (the crux):} ``Write down the answer to (B). Now write the answer to (C) ---
        for Sage, then for Ranger. How many answers did each question have?''
  \item \textbf{2d:} ``Find the shortest wait on the board. Now the longest. Is 25 either one?''
\end{itemize}
While circulating, pick the group work to display: one board with 1c underlined well, and one
group's 2b answer that is \emph{almost} right (many-rows = statistical) to open the debrief with.
\end{multicols}
\end{skillbox}

\boxguard
\begin{skillbox}[Debrief --- Formalize (what goes on the board, in a second color)]{sky}
Work through the displayed student work in order and \textbf{attach each formal name to something
a group already wrote}. Students fill their QuickNotes line by line as you go.
\begin{multicols}{2}
\begin{enumerate}[leftmargin=1.3em, nosep, itemsep=3pt]
  \item Warm-up list of ``things you'd need to know'' $\to$ \emph{context}, and the
        \emph{W questions}. Circle their own words, then write the term beside them.
  \item ``One row is one dog'' $\to$ \emph{individual}. ``Column heading'' $\to$ \emph{variable};
        ``what's written underneath'' $\to$ \emph{value}. Trace a row and a column with a finger
        on the projected board.
  \item 1d, ``the column would be worthless'' $\to$ \emph{variability} --- and say plainly:
        \textbf{this is why the course exists}.
  \item 2b $\to$ the two checks. Write them as the students said them: \emph{do the answers
        differ from one individual to the next?} and \emph{must I collect data?}
  \item Name \emph{statistics} and \emph{data} last, once every part of the definition is
        already on the board.
\end{enumerate}
\columnbreak
\textbf{QuickNotes.} Half a page, filled together. The mini table on the left is the same board
in miniature --- point at it while naming variable / value / individual. The right-hand bullets
are their sentences, tightened. \textbf{It is a summary of what they found, not a new lecture};
if you are talking for more than 90 seconds without pointing at student work, stop.

\textbf{Why this order.} Context first (they already felt its absence in the warm-up), structure
second, variability third, and the statistical-question test last --- because the test
\emph{depends} on variability. Naming the test first would make it a rule to memorize instead of
a consequence.

\textbf{Pose cold:} \emph{``Is `What is the most common breed on this board?' statistical?''}
Breed is not a number, and the answer is one word --- but it varies across dogs and needs the
data. This is the special case the activity does not carry, and it sets up the Application.
\end{multicols}
\end{skillbox}

\boxguard
\begin{skillbox}[Application --- \emph{The Gym Sign-Up} (worked together, 6 min)]{sky}
The first place the just-named vocabulary is \emph{used}. Part (a) is the identification move on
a fresh context; part (b) asks for a statistical question \emph{plus} what varies --- insist on
the second half, because it is the AP habit. \textbf{Part (c) is the real test:} ZIP code is a
label, not a number, and students who have quietly decided ``statistical = numeric'' will say no.
Run both checks aloud on it.

\textbf{Students hold the pen, you hold the questions:} \emph{``Who is one individual here ---
the gym, or a member?''} \quad \emph{``Say your question again, but tell me what differs from
member to member.''} \quad \emph{``Does check~1 say the answers have to be numbers, or just that
they have to differ?''}

\emph{No \texttt{work} block in this lesson:} Topic 1.1 has no computation. The blank and the key
stay page-for-page through \texttt{\textbackslash answerspace} instead.
\end{skillbox}

\boxguard
\begin{skillbox}[Check Your Understanding --- in-class practice, \emph{not scored}]{redbox}
Five exercises spanning Topic 1.1, in new contexts: \textbf{1} identification (turtles);
\textbf{2} write-your-own with what-varies; \textbf{3} the \emph{deliberate contrast pair} --- same
pond, one statistical question and one not; \textbf{4} the bare-claim context critique;
\textbf{5} the AP-style multiple-choice item. Pairs for 1--3, solo for 4--5.

\textbf{They carry no point value} --- the cover's score column reads \textbf{NA}. Do not collect
them for a grade. \textbf{Item~5 is the formative check:} spot-check it on the way out and sort
into three piles --- chose (C) with a variation-based reason (ready); chose (C) with a vague
reason (ready, needs the language); chose (A) or (B) (the misconception survived). If more than a
third land in the last pile, reopen the two checks at the start of Lesson~1.1 rather than moving
on. Eight minutes will not cover five items for everyone: \textbf{priority is 1, 3, and 5};
items~2 and~4 are the early-finisher bank. The graded practice is the homework, below.
\end{skillbox}

\boxguard
\begin{skillbox}[Homework --- AP-style practice (assigned, scored)]{goldbox}
Four AP-style multiple-choice items (five options each, in context) on the adoption board and a
new veterinary context, then one multi-part free-response set on a school-newspaper survey:
identify individuals and variables, write a statistical question with what varies, critique a
bare claim, and explain why a group summary does not predict one individual. The optional
extension is the headline hunt. \textbf{No spiral item --- this is the course's first lesson.}

\textbf{Scoring:} the multiple choice is right/wrong; score the free response the way the AP exam
does --- \emph{in context or it does not count}. A part (b) that names a question but not what
varies is incomplete; a part (d) that restates the claim without mentioning that hours differ
across seniors does not earn the point. Due at the start of Lesson~1.1.
\end{skillbox}

\boxguard
\begin{skillbox}[Watch For (while circulating)]{redbox}
\begin{itemize}[leftmargin=1.2em, nosep]
  \item \textbf{Value named where a variable belongs} (1b, CYU~1, HW~1): ``Beagle'' for ``Breed.''
        Probe: \emph{``Is that the heading, or something written under it?''}
  \item \textbf{``Many rows'' equated with ``statistical''} (2b --- the crux, HW~2): counting the
        dogs needs the whole board and is still not statistical. Probe: \emph{``How many different
        answers does that question have?''}
  \item \textbf{``Has a number in it'' equated with ``statistical''} (CYU~5 distractor B).
        Probe: \emph{``Whose hours? One person's, or many people's?''}
  \item \textbf{``Statistical'' equated with ``numeric''} (Application~c). Probe: \emph{``Run
        check~1. Do members live in different ZIP codes?''}
  \item \textbf{Summary read as a prediction} (2d, HW~d): ``25 days'' taken as a promise.
        Probe: \emph{``Show me a dog on the board that waited 25 days.''}
\end{itemize}
Cold-call prompts: \emph{``Give me one individual and one variable from your board.''} \quad
\emph{``Say a question that is NOT statistical, and tell me which check it fails.''}
\end{skillbox}

\boxguard
\begin{skillbox}[Close \& Preview]{goldbox}
\textbf{Assign the homework} (due next class), then close on what changed rather than on the
assignment: \emph{``This morning seven numbers meant nothing to you. You now know exactly what
you were missing, and you have a two-question test that tells you whether a question is even
worth collecting data for.''} \textbf{Preview:} Lesson~1.1, \emph{The Language of Variation} ---
we split variables into categorical and quantitative, the first decision that determines every
graph and summary for the rest of the course. Note that \emph{population} and \emph{sample} are
deliberately \textbf{not} introduced today; they belong to Unit~3, and naming them now competes
with the individual/variable distinction this lesson is built on.
\end{skillbox}

% ── Teacher notes — one per component, in packet order ────────────────────────
\begin{teachernote}[Warm-Up]
Deliberately unanswerable at first glance. Give three minutes and resist rescuing --- the
discomfort is the lesson. Expect some students to guess ``test scores'' and defend it; validate
the reasoning while pointing out that a guess is not information. Debrief only item~3 aloud and
leave it on the board all period; it is the same discrimination the activity's crux tests. If it
runs long, cut item~2 and pick the ``what would you need to know'' list up during the debrief
instead --- but never cut item~1, which is the whole seed.
\end{teachernote}

\begin{teachernote}[Experience \& Formalize]
\textbf{Activity (20 min).} Roughly 8 minutes on scenario~1, 12 on scenario~2 --- scenario~2
carries the crux and needs the time. Groups that finish scenario~1 early should be pushed to
justify 1d out loud rather than moving on. \textbf{The single discipline that makes this lesson
work: do not name anything while circulating.} If you say ``variable'' at a table, that table
stops discovering. Ask where they see it instead.

\textbf{Debrief (12 min).} The must-land moments are (2) heading-vs-entry and (4) the two checks;
everything else can be compressed. Open on the \emph{almost}-right 2b answer you collected --- the
group that decided many-rows means statistical --- because correcting it in public, gently, is
what makes the crux stick. If time is short, cut the cold-call breed question and pick it up as
the Application's part (c), which tests the same thing.

\textbf{Application (6 min).} Part (c) is the point of the whole box; if you only get through
(a) and (c), that is the right trade. Students hold the pen.

\textbf{Check Your Understanding (8 min).} Not scored --- spot-check rather than collecting the
page. Priority items are 1, 3, and 5; 2 and 4 are the early-finisher bank. Item~3 is the
deliberate contrast pair and item~5 is the formative check that decides how Lesson~1.1 opens.
\end{teachernote}

\begin{teachernote}[Homework]
The adoption board is reused on purpose --- the context is already loaded, so the thinking stays
on the distinctions rather than on re-reading a table. MC item~1 targets variable-vs-value, item~2
the crux, item~3 identification in a fresh context, and item~4 variability as the reason the
subject exists. \textbf{Item~4 is the one worth reading closely}: it is EK VAR-1.A.1 in its most
transferable form and recurs on the AP exam as ``interpret in context.'' On the free response,
hold the line on context from day one --- part (d) is where students first learn that an answer
without the situation in it earns nothing. The extension pairs well with a bulletin-board display
of headlines students bring in.
\end{teachernote}
\end{document}
